% =====================================================================
% 1. Introducción
% =====================================================================
\section{Introducción}
\label{sec:intro}

\subsection{Contexto y objetivo}

La UPME evalúa la factibilidad de granjas solares en distintas regiones del
país y quiere priorizar los municipios designados como PDET (Programas de
Desarrollo con Enfoque Territorial), un foco de la política de desarrollo
posconflicto~\cite{pdet}. A mayor área de techos en un municipio, mayor
potencial para captar energía con paneles instalados sobre ellos. El
enunciado del proyecto~\cite{enunciado} pide una metodología reproducible
para \textbf{contar las edificaciones y sumar el área de sus techos en cada
municipio PDET}, comparando al menos dos fuentes abiertas de huellas de
edificios y usando una solución NoSQL.

El listado oficial (\campo{MunicipiosPDET.xlsx}) tiene 170 municipios en 16
subregiones PDET y 19 departamentos. Los límites municipales salen del Marco
Geoestadístico Nacional (MGN) del DANE~\cite{mgn}.

\subsection{Alcance}

\begin{itemize}
  \item \textbf{Territorio:} los 170 municipios PDET, con sus límites del MGN
        2025 (389\,182\,km\textsuperscript{2}).
  \item \textbf{Fuentes de edificios:} Google Open Buildings v3~\cite{google}
        y Microsoft Building Footprints~\cite{microsoft}, dos de las tres que
        propone el enunciado (la tercera es GlobalBuildingAtlas). La
        justificación está en la tabla~\ref{tab:decisiones}.
  \item \textbf{Métricas por municipio y fuente:} número de edificios, área
        total de techos, área media, densidad (edificios por
        km\textsuperscript{2}) y porcentaje del municipio cubierto por techos,
        con un análisis de sensibilidad al umbral de confianza de Google.
  \item \textbf{Fuera de alcance:} el diseño de los pilotos, la irradiación
        por municipio y la demanda eléctrica. La sección~\ref{sec:potencial}
        da solo un orden de magnitud de la potencia instalable.
\end{itemize}

El área de techo de cada edificio se aproxima con el área de su huella en
planta, que es lo que entregan las dos fuentes. Para cada municipio $m$ y
fuente $f$:

\begin{equation}
  A_{m,f} = \sum_{e \in E_f,\; \mathrm{centroide}(e) \in m} \mathrm{area}_{9377}(e)
  \label{eq:area}
\end{equation}

donde $E_f$ es el conjunto de edificios de la fuente $f$ y
$\mathrm{area}_{9377}(e)$ es el área de la huella calculada en la proyección
plana oficial de Colombia (EPSG:9377, en metros).

\subsection{Cómo se construyó y cómo leer este informe}

El proyecto se desarrolló en cuatro entregas semanales, cada una con su
código en el repositorio y su informe: diseño del esquema (semana 1), límites
PDET (semana 2), huellas de edificios (semana 3) y análisis (semana 4). Este
informe las integra en un solo documento:

\begin{itemize}
  \item Secciones~\ref{sec:arquitectura} y~\ref{sec:modelo}: arquitectura,
        decisiones tecnológicas y modelo de datos.
  \item Secciones~\ref{sec:municipios} a~\ref{sec:eda}: adquisición,
        integridad y carga de los datos, y su análisis exploratorio.
  \item Sección~\ref{sec:metodologia}: cómo se calcularon las métricas.
  \item Secciones~\ref{sec:resultados} y~\ref{sec:discusion}: resultados,
        discusión y recomendaciones.
  \item Sección~\ref{sec:reproduccion}: cómo reproducir todo.
  \item Anexo~\ref{sec:anexo}: capturas de evidencia de cada etapa.
\end{itemize}
